\textbf{Phần 1. Trắc nghiệm lựa chọn}
\setcounterpageref{ex}{0}

\begin{ex} %[5D4N20]
Kết quả của phép tính $5{,}8 - 3{,}9$ là
\choice
{\True $1{,}9$}
{$2{,}9$}
{$1{,}1$}
{$2{,}1$}
\loigiai{
Ta có: $5{,}8 - 3{,}9 = 1{,}9$.\\
\textbf{Đáp án của câu hỏi là A}
}
\end{ex}

\begin{ex} %[5D4N20]
Kết quả của phép tính $63{,}49 - 1{,}8$ là
\choice
{$62{,}41$}
{\True $61{,}69$}
{$61{,}41$}
{$45{,}49$}
\loigiai{
Đặt tính thẳng cột: $63{,}49 - 1{,}80 = 61{,}69$.\\
\textbf{Đáp án của câu hỏi là B}
}
\end{ex}

\begin{ex} %[5D4H20]
Kết quả của phép tính $17{,}96 - 8{,}5$ là
\choice
{\True $9{,}46$}
{$9{,}91$}
{$9{,}41$}
{$17{,}11$}
\loigiai{
Ta có: $17{,}96 - 8{,}50 = 9{,}46$.\\
\textbf{Đáp án của câu hỏi là A}
}
\end{ex}

\begin{ex} %[5D4H20]
Kết quả của phép tính $25{,}9 - 13{,}84$ là
\choice
{$12{,}14$}
{\True $12{,}06$}
{$12{,}16$}
{$12{,}86$}
\loigiai{
Ta có: $25{,}90 - 13{,}84 = 12{,}06$.\\
\textbf{Đáp án của câu hỏi là B}
}
\end{ex}

\begin{ex} %[5D4H20]
Kết quả của phép tính $7{,}6 - 1{,}51$ là
\choice
{$6{,}19$}
{$6{,}11$}
{$5{,}09$}
{\True $6{,}09$}
\loigiai{
Ta có: $7{,}60 - 1{,}51 = 6{,}09$.\\
\textbf{Đáp án của câu hỏi là D}
}
\end{ex}

\begin{ex} %[5D4H20]
Kết quả của phép tính $21{,}4 - 6$ là
\choice
{$20{,}8$}
{$14{,}6$}
{\True $15{,}4$}
{$15{,}6$}
\loigiai{
Trừ phần nguyên: $21 - 6 = 15$, giữ nguyên phần thập phân: $15{,}4$.\\
\textbf{Đáp án của câu hỏi là C}
}
\end{ex}

\begin{ex} %[5D4H20]
Kết quả của phép tính $9 - 3{,}5$ là
\choice
{$6{,}5$}
{$5{,}0$}
{$6{,}0$}
{\True $5{,}5$}
\loigiai{
Ta có: $9{,}0 - 3{,}5 = 5{,}5$.\\
\textbf{Đáp án của câu hỏi là D}
}
\end{ex}

\begin{ex} %[5D4H20]
Tìm số thích hợp điền vào dấu hỏi chấm: $? + 8{,}16 = 17{,}5$.
\choice
{$25{,}66$}
{$9{,}46$}
{\True $9{,}34$}
{$8{,}34$}
\loigiai{
Ta có: $? = 17{,}5 - 8{,}16 = 17{,}50 - 8{,}16 = 9{,}34$.\\
\textbf{Đáp án của câu hỏi là C}
}
\end{ex}

\begin{ex} %[5D4H20]
Tìm số thích hợp điền vào dấu hỏi chấm: $? - 6{,}17 = 11{,}83$.
\choice
{$5{,}66$}
{$17{,}00$}
{$18{,}1$}
{\True $18$}
\loigiai{
Ta có: $? = 11{,}83 + 6{,}17 = 18{,}00 = 18$.\\
\textbf{Đáp án của câu hỏi là D}
}
\end{ex}

\begin{ex} %[5D4V20]
Mai làm được $2{,}15\text{ l}$ nước mơ, Mi làm được $1{,}7\text{ l}$ nước dâu. Hỏi lượng nước mơ của Mai nhiều hơn lượng nước dâu của Mi bao nhiêu lít?
\choice
{$0{,}35\text{ l}$}
{\True $0{,}45\text{ l}$}
{$1{,}45\text{ l}$}
{$0{,}045\text{ l}$}
\loigiai{
Nước mơ nhiều hơn nước dâu là: $2{,}15 - 1{,}7 = 0{,}45\text{ (l)}$.\\
\textbf{Đáp án của câu hỏi là B}
}
\end{ex}

\begin{ex} %[5D4V20]
Một chiếc cọc được sơn hai màu xanh và đỏ. Phần màu xanh dài $14{,}2\text{ dm}$ và dài hơn đoạn màu đỏ $1{,}8\text{ dm}$. Tổng chiều dài của chiếc cọc đó là
\choice
{\True $26{,}6\text{ dm}$}
{$12{,}4\text{ dm}$}
{$16{,}0\text{ dm}$}
{$28{,}4\text{ dm}$}
\loigiai{
Đoạn màu đỏ dài: $14{,}2 - 1{,}8 = 12{,}4\text{ (dm)}$.\\
Chiều dài cả chiếc cọc là: $14{,}2 + 12{,}4 = 26{,}6\text{ (dm)}$.\\
\textbf{Đáp án của câu hỏi là A}
}
\end{ex}

\begin{ex} %[5D4V20]
Cả ba chú rô-bốt A, B, C có tổng cân nặng là $8\text{ kg}$. Rô-bốt A và B cân nặng $4{,}7\text{ kg}$, rô-bốt B và C cân nặng $5{,}5\text{ kg}$. Cân nặng của rô-bốt A là
\choice
{$2{,}2\text{ kg}$}
{$3{,}3\text{ kg}$}
{\True $2{,}5\text{ kg}$}
{$3{,}5\text{ kg}$}
\loigiai{
Cân nặng của rô-bốt A là: $8 - 5{,}5 = 2{,}5\text{ (kg)}$.\\
\textbf{Đáp án của câu hỏi là C}
}
\end{ex}

\textbf{Phần 2. Trắc nghiệm đúng sai}
\setcounterpageref{ex}{0}

\begin{ex} %[5D4H20]
Xét tính đúng sai của các phép tính sau:
\choiceTF
{Phép trừ $6{,}15 - 2{,}7 = 4{,}8$.}
{\True Phép trừ $4{,}38 - 1{,}56 = 2{,}82$.}
{Phép trừ $85{,}9 - 5{,}06 = 35{,}3$.}
{\True Phép trừ $4{,}43 - 4{,}16 = 0{,}27$.}
\loigiai{
\begin{itemize}
    \item $6{,}15 - 2{,}70 = 3{,}45 \neq 4{,}8$.
    \item $4{,}38 - 1{,}56 = 2{,}82$.
    \item $85{,}90 - 5{,}06 = 80{,}84 \neq 35{,}3$ (sai do đặt tính lệch hàng).
    \item $4{,}43 - 4{,}16 = 0{,}27$.
\end{itemize}
\begin{itemize}
    \item a) \textbf{Sai}: Kết quả là $3{,}45$.
    \item b) \textbf{Đúng}.
    \item c) \textbf{Sai}: Kết quả là $80{,}84$.
    \item d) \textbf{Đúng}.
\end{itemize}
}
\end{ex}

\begin{ex} %[5D4V20]
Cho ba chú rô-bốt A, B, C có tổng cân nặng là $8\text{ kg}$. Biết rô-bốt A và B nặng $4{,}7\text{ kg}$; rô-bốt B và C nặng $5{,}5\text{ kg}$. Xét tính đúng sai của các khẳng định sau:
\choiceTF
{\True Rô-bốt C cân nặng $3{,}3\text{ kg}$.}
{\True Rô-bốt B cân nặng $2{,}2\text{ kg}$.}
{Rô-bốt A nặng hơn rô-bốt C là $0{,}8\text{ kg}$.}
{\True Trong ba chú rô-bốt, rô-bốt B là nhẹ nhất.}
\loigiai{
\begin{itemize}
    \item Cân nặng của rô-bốt C: $8 - 4{,}7 = 3{,}3\text{ kg}$.
    \item Cân nặng của rô-bốt A: $8 - 5{,}5 = 2{,}5\text{ kg}$.
    \item Cân nặng của rô-bốt B: $4{,}7 - 2{,}5 = 2{,}2\text{ kg}$.
    \item Rô-bốt C nặng hơn rô-bốt A: $3{,}3 - 2{,}5 = 0{,}8\text{ kg}$ (Rô-bốt A nhẹ hơn rô-bốt C).
    \item So sánh: $2{,}2 < 2{,}5 < 3{,}3 \Rightarrow$ Rô-bốt B nhẹ nhất.
\end{itemize}
\begin{itemize}
    \item a) \textbf{Đúng}.
    \item b) \textbf{Đúng}.
    \item c) \textbf{Sai}: Rô-bốt A nhẹ hơn rô-bốt C.
    \item d) \textbf{Đúng}.
\end{itemize}
}
\end{ex}

\textbf{Phần 3. Trả lời ngắn}
\setcounterpageref{ex}{0}

\begin{ex} %[5D4H20]
Tính hiệu của hai số $2{,}53$ và $1{,}62$. (Nhập kết quả dưới dạng số thập phân).
\shortans{0,91}
\loigiai{
\textbf{Đáp số:} $0{,}91$\\
$2{,}53 - 1{,}62 = 0{,}91$.
}
\end{ex}

\begin{ex} %[5D4H20]
Tính giá trị của biểu thức: $25{,}9 - 13{,}84$. (Nhập kết quả dưới dạng số thập phân).
\shortans{12,06}
\loigiai{
\textbf{Đáp số:} $12{,}06$\\
$25{,}90 - 13{,}84 = 12{,}06$.
}
\end{ex}

\begin{ex} %[5D4H20]
Tìm số thích hợp điền vào dấu hỏi chấm: $8{,}9 + ? = 28{,}501$. (Nhập kết quả dưới dạng số thập phân).
\shortans{19,601}
\loigiai{
\textbf{Đáp số:} $19{,}601$\\
$? = 28{,}501 - 8{,}9 = 19{,}601$.
}
\end{ex}

\begin{ex} %[5D4V20]
Thành tích nhảy xa năm nay của Nam là $4{,}43\text{ m}$, năm ngoái là $4{,}16\text{ m}$. Hỏi thành tích năm nay của Nam hơn năm ngoái bao nhiêu mét? (Nhập kết quả dưới dạng số thập phân).
\shortans{0,27}
\loigiai{
\textbf{Đáp số:} $0{,}27$\\
$4{,}43 - 4{,}16 = 0{,}27\text{ (m)}$.
}
\end{ex}

\begin{ex} %[5D4V20]
Một chiếc cọc có phần màu xanh dài $14{,}2\text{ dm}$ và phần màu đỏ ngắn hơn phần màu xanh $1{,}8\text{ dm}$. Hỏi toàn bộ chiếc cọc dài bao nhiêu đề-xi-mét? (Nhập kết quả dưới dạng số thập phân).
\shortans{26,6}
\loigiai{
\textbf{Đáp số:} $26{,}6$\\
Phần màu đỏ: $14{,}2 - 1{,}8 = 12{,}4\text{ dm}$. Chiếc cọc dài: $14{,}2 + 12{,}4 = 26{,}6\text{ dm}$.
}
\end{ex}

\begin{ex} %[5D4V20]
Cả ba rô-bốt A, B, C cân nặng $8\text{ kg}$, rô-bốt A và B cân nặng $4{,}7\text{ kg}$, rô-bốt B và C cân nặng $5{,}5\text{ kg}$. Hỏi riêng chú rô-bốt B cân nặng bao nhiêu ki-lô-gam? (Nhập kết quả dưới dạng số thập phân).
\shortans{2,2}
\loigiai{
\textbf{Đáp số:} $2{,}2$\\
Rô-bốt A: $8 - 5{,}5 = 2{,}5\text{ kg}$. Rô-bốt B: $4{,}7 - 2{,}5 = 2{,}2\text{ kg}$.
}
\end{ex}
